\documentclass[11pt,a4paper]{article}

\def\ncourse {Introduction to Functional Analysis and Fourier Analysis}
\def\ntype {problems}
\def\ncoursehead {Functional Analysis 104}

\makeatletter

% packages
\usepackage{amssymb,amsfonts,amsmath,calc,tikz,pgfplots,geometry,mathtools}
\usepackage{color}   % May be necessary if you want to color links
\usepackage[svgnames]{xcolor}
\usepackage[hidelinks]{hyperref}
\usepackage{forest}
\usepackage{commath} % ew
\usepackage{esvect} % \vv makes vectors
\usepackage{centernot} % for the comparison
\usepackage{esdiff}
\usepackage{esint}
\usepackage{amsthm}
\usepackage{fancyhdr}
\usepackage{bm}
\usepackage{witharrows}
\usepackage{bookmark}
\usepackage{tikz-cd}
\usepackage{quiver}
\usepackage{bbm}
\usepackage{textcomp}
\usepackage{float}
\usepackage{gensymb}
\usepackage{cleveref}
\usepackage{environ}
\usepackage{comment}
\usepackage{cancel}
\usepackage{xparse}

% Inevitable cs
\usepackage{listings}
\usepackage[]{algorithm2e}

% tikz libraries
\usetikzlibrary{positioning}
\usetikzlibrary{matrix}
\usetikzlibrary{arrows}
\usetikzlibrary{bending}
\usetikzlibrary{arrows.meta}
\usetikzlibrary{decorations.markings}
\usetikzlibrary{decorations.pathreplacing}
\usetikzlibrary{decorations.markings}
\usetikzlibrary{patterns}
\usetikzlibrary{backgrounds}

% Page style setup
\pagestyle{fancy}
\fancyhfoffset{0pt}
\renewcommand{\headwidth}{\textwidth}
\geometry{margin=1in}
\pgfplotsset{compat=1.18}
\setlength{\headheight}{14.6pt}
\addtolength{\topmargin}{-1.6pt}
\hypersetup{
    colorlinks=false,
    linktoc=section,
    linkcolor=black,
}

%% maketitle setup
\ifx \nauthor\undefined
  \def\nauthor{Yeheli Fomberg}
\else
\fi

\ifx \nemail\undefined
  \def\nemail{\href{mailto:yp@campus.technion.ac.il}
  {\texttt{<yp@campus.technion.ac.il>}}}
\else
\fi

\ifx \ncoursehead \undefined
  \def\ncoursehead{\ncourse}
\fi

\lhead{\emph{\nouppercase{\leftmark}}}
\ifx \nextra \undefined
  \rhead{
    \ifnum\thepage=1
    \else
      \ncoursehead
    \fi}
\else
  \rhead{
    \ifnum\thepage=1
    \else
      \ncoursehead \ (\nextra)
    \fi}
\fi

\def\notes{notes}
\def\problems{problems}

\ifx \ntype \undefined
  \def\ntype{notes}
\else
\fi

\ifx \ntype \notes
  \let\@real@maketitle\maketitle
  \renewcommand{\maketitle}{\@real@maketitle\begin{center}
  \begin{minipage}[c]{0.9\textwidth}\centering\footnotesize
  These notes are not endorsed by the lecturers.
  I have revised them outside lectures to incorporate supplementary
  explanations, clarifications, and material for fun.
  While I have strived for accuracy, any errors or misinterpretations 
  are most likely mine.
  \end{minipage}\end{center}}
\else
  \ifx \ntype \problems
    \let\@real@maketitle\maketitle
    \renewcommand{\maketitle}{\@real@maketitle\begin{center}
    \begin{minipage}[c]{0.9\textwidth}\centering\footnotesize
    These problems are mostly based on the problem sets I was given
    when taking the course.
    
    While I have strived for accuracy and completeness,
    some of the problems would not have been here without the help of
    friends.
    \end{minipage}\end{center}}
  \else
  \fi
\fi


% theorem environments
\theoremstyle{definition}
\newtheorem{definition}{Definition}[section]
\newtheorem{remark}{Remark}[section]
\newtheorem{example}{Example}[section]
\newtheorem{exercise}{Exercise}[section]
\newtheorem{paradox}{Paradox}[section]
\newtheorem{entry}{Entry}[section]
\newtheorem*{solution}{Solution}
\newtheorem*{question}{Question}
\newtheorem*{answer}{Answer}
\newtheorem*{problem}{Problem}
\theoremstyle{plain}
\newtheorem{theorem}{Theorem}[section]
\newtheorem{proposition}[theorem]{Proposition}
\newtheorem{lemma}[theorem]{Lemma}
\newtheorem{corollary}[theorem]{Corollary}
\newenvironment{proofof}[1]{%
    \renewcommand{\proofname}{Proof of #1}%
    \begin{proof}
}{%
    \end{proof}
}
% tikz customization
\pgfarrowsdeclarecombine{twolatex'}{twolatex'}{latex'}{latex'}{latex'}{latex'}
\tikzset{->/.style = {decoration={markings,
                                  mark=at position 0.999
                                  with {\arrow[scale=2]{latex'}}},
                      postaction={decorate}}}
\tikzset{<-/.style = {decoration={markings,
                                  mark=at position 0.001 with {\arrowreversed[scale=2]{latex'}}},
                      postaction={decorate}}}
\tikzset{-->/.style = {decoration={markings,
                                  mark=at position 0.999
                                  with {\arrow[scale=2]{latex'[sep=0pt -2.5]}}},
                      postaction={decorate}}}
\tikzset{<--/.style = {decoration={markings,
                                  mark=at position 0.001
                                  with {\arrowreversed[scale=2]{latex'[sep=0pt -2.5]}}},
                      postaction={decorate}}}
\tikzset{<->/.style = {decoration={markings,
                                  mark=at position 0.001 with {\arrowreversed[scale=2]{latex'}},
                                  mark=at position 0.999 with {\arrow[scale=2]{latex'}},},
                      postaction={decorate}}}
\tikzset{->-/.style = {decoration={markings,
    mark=at position 0.50+0.225/#1 with {\arrow[scale=2]{latex'}}},
                      postaction={decorate}}}
\tikzset{-<-/.style = {decoration={markings,
                                   mark=at position 0.50-0.225/#1 with {\arrowreversed[scale=2]{latex'}}},
                       postaction={decorate}}}
\tikzset{->>/.style = {decoration={markings,
                                  mark=at position 1 with {\arrow[scale=2]{latex'}}},
                      postaction={decorate}}}
\tikzset{<<-/.style = {decoration={markings,
                                  mark=at position 0 with {\arrowreversed[scale=2]{twolatex'}}},
                      postaction={decorate}}}
\tikzset{<<->>/.style = {decoration={markings,
                                   mark=at position 0 with {\arrowreversed[scale=2]{twolatex'}},
                                   mark=at position 1 with {\arrow[scale=2]{twolatex'}}},
                       postaction={decorate}}}
\tikzset{->>-/.style = {decoration={markings,
                                   mark=at position 0.50+0.450/#1 with {\arrow[scale=2]{twolatex'}}},
                       postaction={decorate}}}
\tikzset{-<<-/.style = {decoration={markings,
                                   mark=at position 0.50-0.450/#1 with {\arrowreversed[scale=2]{twolatex'}}},
                       postaction={decorate}}}

\pgfarrowsdeclare{biggertip}{biggertip}{%
  \setlength{\arrowsize}{1pt}
  \addtolength{\arrowsize}{.1\pgflinewidth}
  \pgfarrowsrightextend{0}
  \pgfarrowsleftextend{-5\arrowsize}
}{%
  \setlength{\arrowsize}{1pt}
  \addtolength{\arrowsize}{.1\pgflinewidth}
  \pgfpathmoveto{\pgfpoint{-5\arrowsize}{4\arrowsize}}
  \pgfpathlineto{\pgfpointorigin}
  \pgfpathlineto{\pgfpoint{-5\arrowsize}{-4\arrowsize}}
  \pgfusepathqstroke
}
\tikzset{
	EdgeStyle/.style = {>=biggertip}
}

\tikzset{circ/.style = {fill, circle, inner sep = 0, minimum size = 3}}
\tikzset{scirc/.style = {fill, circle, inner sep = 0, minimum size = 1.5}}
\tikzset{mstate/.style={circle, draw, black, text=black, minimum width=0.7cm}}

\tikzset{eqpic/.style={baseline={([yshift=-.5ex]current bounding box.center)}}}

\definecolor{mblue}{rgb}{0.2, 0.3, 0.8}
\definecolor{morange}{rgb}{1, 0.5, 0}
\definecolor{mgreen}{rgb}{0, 0.4, 0.2}
\definecolor{mred}{rgb}{0.5, 0, 0}

% topology
\DeclareMathOperator{\dfr}{dfr} % deformation retract
\newcommand{\Cells}{\text{Cells}}
\newcommand{\RP}{\mathbb{RP}}
\newcommand{\CP}{\mathbb{CP}}

% order theory
\newcommand{\join}{\vee}
\newcommand{\meet}{\wedge}

% algebraic geometry
\DeclareMathOperator{\Rad}{Rad} % Radical of modules
\DeclareMathOperator{\Spec}{Spec}
% \renewcommand{\div}{\mathrm{div}} % divisors
\DeclareMathOperator{\Mer}{Mer} % Meromorphic functions
\DeclareMathOperator{\Exc}{Exc} % Exceptional divisor
\DeclareMathOperator{\Td}{Td} % Todd classes

% algebra

\DeclareMathOperator{\trdeg}{tr.deg.}
\DeclareMathOperator{\odd}{odd}
\DeclareMathOperator{\even}{even}
\DeclareMathOperator{\lcm}{lcm}
\DeclareMathOperator{\ord}{ord}
\DeclareMathOperator{\codim}{codim}
\DeclareMathOperator{\Ann}{Ann}
\DeclareMathOperator{\Ext}{Ext}
\DeclareMathOperator{\Sym}{Sym}
\DeclareMathOperator{\Out}{Out}
\DeclareMathOperator{\Aut}{Aut}
\DeclareMathOperator{\End}{End}
\DeclareMathOperator{\Inn}{Inn}
\DeclareMathOperator{\Mat}{Mat}
\DeclareMathOperator{\Fix}{Fix}
\DeclareMathOperator{\std}{std}
\DeclareMathOperator{\sgn}{sgn}
\DeclareMathOperator{\adj}{adj}
\DeclareMathOperator{\id}{id}
\DeclareMathOperator{\op}{op}
\DeclareMathOperator{\tr}{tr}
\DeclareMathOperator{\diag}{diag}
\DeclareMathOperator{\rank}{rank}
\DeclareMathOperator{\rk}{rk}
\DeclareMathOperator{\coker}{coker}
\DeclareMathOperator{\Mult}{Mult} % multiplier algebra
\DeclareMathOperator{\Frac}{Frac} % fraction field
\DeclareMathOperator{\Rat}{Rat}
\DeclareMathOperator{\Gal}{Gal}
\newcommand{\ioplus}{\overset{\text{int}}{\oplus}} % interior direct sum
\newcommand{\GG}{\mathcal G} % Galois correspondence
\newcommand{\FF}{\mathcal F} % Galois correspondence
\newcommand{\cupp}{\smile} % cup product
\newcommand{\capp}{\frown} % cap product
\newcommand{\actson}{\curvearrowright}
\newcommand{\bra}{\left\langle}
\newcommand{\ket}{\right\rangle}
\newcommand{\idealin}{\triangleleft\,}
\newcommand{\normalin}{\triangleleft\,}
\newcommand{\ip}[2]{\left\langle #1, #2 \right\rangle}
\newcommand{\bigslant}[2]
{{\raisebox{.2em}{$#1$}\left/\raisebox{-.2em}{$#2$}\right.}}
\newcommand{\properideal}{%
  \mathrel{\ooalign{$\lneq$\cr\raise.22ex\hbox{$\lhd$}\cr}}}

% categories
\newcommand{\cat}[1]{\mathbf{#1}}
\newcommand{\Ch}{\cat{Ch}} % Chain complexes   
\newcommand{\Set}{\cat{Set}}
\newcommand{\Grp}{\cat{Grp}}
\newcommand{\Ab}{\cat{Ab}}
\newcommand{\Ring}{\cat{Ring}}
\newcommand{\CRing}{\cat{CRing}}
\newcommand{\Top}{\cat{Top}}
\newcommand{\Vect}{\cat{Vect}}
\newcommand{\cmod}{\cat{mod}}

%% matrix groups
\newcommand{\GL}{\mathrm{GL}}
\newcommand{\Or}{\mathrm{O}}
\newcommand{\PGL}{\mathrm{PGL}}
\newcommand{\PSL}{\mathrm{PSL}}
\newcommand{\PSO}{\mathrm{PSO}}
\newcommand{\PSU}{\mathrm{PSU}}
\newcommand{\SL}{\mathrm{SL}}
\newcommand{\msl}{\mathrm{sl}}
\newcommand{\SO}{\mathrm{SO}}
\newcommand{\Spin}{\mathrm{Spin}}
\newcommand{\Sp}{\mathrm{Sp}}
\newcommand{\SU}{\mathrm{SU}}
\newcommand{\U}{\mathrm{U}}
\let\char\relax
\newcommand{\char}{\text{char}}

% analysis
\renewcommand{\Re}{\mathrm{Re}}
\renewcommand{\Im}{\mathrm{Im}}
\NewDocumentCommand{\dx}{o}{
  \IfNoValueTF{#1}
    {\dif x}                  % if no optional argument
    {\dif^{#1} \!x}         % if optional argument
}
\NewDocumentCommand{\dy}{o}{
  \IfNoValueTF{#1}
    {\dif y}                  % if no optional argument
    {\dif^{#1} \!y}         % if optional argument
}
\newcommand{\dt}{\dif t}
\newcommand{\du}{\dif u}
\newcommand{\dv}{\dif v}
\newcommand{\dz}{\dif z}
\newcommand{\ds}{\dif s}
\newcommand{\dw}{\dif w}
\newcommand{\dr}{\dif r}
\newcommand{\dm}{\dif m}
\newcommand{\df}{\dif f}
\newcommand{\dV}{\dif V}
\newcommand{\dS}{\dif S}
\newcommand{\dA}{\dif \mathrm{A}}
\newcommand{\dmu}{\dif \mu}
\newcommand{\dnu}{\dif \nu}
\newcommand{\dtheta}{\dif \theta}
\newcommand{\domega}{\dif \omega}
\newcommand{\dsigma}{\dif \sigma}
\newcommand{\dtau}{\dif \tau}
\newcommand{\dvarphi}{\dif \varphi}
\renewcommand{\dif}{\mathrm{d}}
\renewcommand{\Dif}{\mathrm{D}}
\renewcommand{\triangle}{\Delta}
\newcommand{\ptriangle}{\partial \triangle}
\DeclareMathOperator{\im}{im}
\DeclareMathOperator{\cis}{cis}
\DeclareMathOperator{\rot}{rot}
\DeclareMathOperator{\vspan}{span}
\DeclareMathOperator{\grad}{grad}
\DeclareMathOperator{\curl}{curl}
\DeclareMathOperator{\esssup}{esssup}
\newcommand{\hodge}{{\mathnormal{\star}}}
\DeclareMathOperator{\range}{range}
\DeclareMathOperator{\Hom}{Hom}
\DeclareMathOperator{\Int}{Int}
\DeclareMathOperator{\Conv}{Conv} % convex hull
\newcommand{\ind}{\mathbbm{1}}
\DeclareMathOperator{\Res}{Res}
\DeclareMathOperator{\graph}{graph}
\DeclareMathOperator{\diam}{diam}
\DeclareMathOperator{\supp}{supp}
\DeclareMathOperator{\Vol}{Vol} % Volume
\DeclareMathOperator{\Area}{Area} % Area
\DeclareMathOperator{\area}{area}
\DeclareMathOperator{\Ind}{Ind} % Index
\newcommand{\length}{\mathrm{length}}

% logic
\DeclareMathOperator{\MOD}{MOD}
\DeclareMathOperator{\Theory}{Theory}
\newcommand{\notimplies}{%
  \mathrel{{\ooalign{\hidewidth$\not\phantom{=}$\hidewidth\cr$\implies$}}}}

% nice
\newcommand{\half}{\frac{1}{2}}
\newcommand{\pair}{\del}
\newcommand{\taking}[1]{\xrightarrow{#1}}
\newcommand{\otaking}[1]{\xleftarrow{#1}}
\newcommand{\inv}{^{-1}}
\newcommand{\ot}{\leftarrow}
\newcommand{\ninfty}{-\infty}
\newcommand{\pinfty}{+\infty}
\newcommand{\floor}[1]{\left\lfloor #1 \right\rfloor}
\newcommand{\ceil}[1]{\left\lceil #1 \right\rceil}
\newcommand{\viff}{\Big\Updownarrow} % vertical iff
\newcommand{\bmid}{\,\middle\vert\,} % big mid

% probability
\newcommand{\Prob}{\mathbf{P}}
\renewcommand{\vec}[1]{\boldsymbol{\mathbf{#1}}}
\DeclareMathOperator{\Bin}{Bin}
\DeclareMathOperator{\Geo}{Geo}
\DeclareMathOperator{\Poi}{Poi}
\DeclareMathOperator{\Exp}{Exp}
\DeclareMathOperator{\Var}{Var} % Variance
\DeclareMathOperator{\Cov}{Cov}

% special letters
\newcommand{\N}{\mathbb{N}}
\newcommand{\Z}{\mathbb{Z}}
\newcommand{\Q}{\mathbb{Q}}
\newcommand{\R}{\mathbb{R}}
\newcommand{\C}{\mathbb{C}}
\newcommand{\F}{\mathbb{F}}
\newcommand{\E}{\mathbf{E}}
\newcommand{\D}{\mathbb{D}}
\renewcommand{\H}{\mathbb{H}}
\newcommand{\ps}{\mathcal{P}}
\newcommand{\M}{\mathcal{M}}
\renewcommand{\L}{\mathcal{L}}
\renewcommand{\O}{\mathcal{O}}
\renewcommand{\U}{\mathcal{U}}
\newcommand{\Omicron}{O}
\newcommand{\powerset}{\mathcal{P}}

% text
\newcommand{\st}{\text{ s.t. }}
\newcommand{\tand}{\quad \text{and} \quad}
\newcommand{\tor}{\quad \text{or} \quad}
\newcommand{\stand}{\text{ and }}
\newcommand{\stor}{\text{ or }}
\renewcommand{\tt}[1]{\textnormal{\textbf{(#1).}}} %tt=theorem title GET RID OF

% title format
\title{\textbf{\ncourse}}

\ifx \ntype \notes
  \author{Notes taken by \nauthor\, \nemail \\ Based on lectures by \nlecturer}
  \date{\nterm\ \nyear}
\else
  \ifx \ntype \problems
    \author{These problems were solved by \nauthor \\ \nemail}
  \else
  \fi
\fi

\makeatother


\begin{document}
\maketitle

\newpage
\tableofcontents
\newpage

\section{Introduction and the Stone--Weierstrass theorem}

\begin{exercise}
  Let $f \in C_{\R}([-1,1])$ and assume that $\int_{-1}^{1} f(x) x^n \dx = 0$,
  for all odd positive integers $n$.
  Prove that $f$ is an even function.
\end{exercise}
\begin{solution}
  First we will show that if $f$ is a continuous real-valued function on $[a,b]$
  such that $\int_{a}^{b} f(x) x^n \dx = 0$ for all $n = 0,1,2,\dots$
  then $f \equiv 0$.
  We will show this by proving that the integral of $f^2$ over $[a,b]$
  is $0$.
  The result will follow because if $f$ is not the zero function,
  from continuity $f^2$ will have an interval with values greater than
  some $\epsilon > 0$, which contradicts $\int f^2 = 0$.
  From linearity of the integral, and what is given, we know that
  \[
    \int_{a}^{b} f(x) p(x) \dx = 0
  \]
  for any polynomial $p(x) \in \R[x]$.
  From the Weierstrass approximation theorem for every $1/n > 0$
  there exists a polynomial $p_n(x) \in C[x]$ so that
  $\norm{f - p}_{\infty} < 1 / n$.
  Notice that we have for every $n \geq 1$ that
  \begin{align*}
    \abs{\int_{a}^{b} f(x)f(x) \dx} =
    \abs{\int_{a}^{b} f(x)f(x) \dx - \int_{a}^{b} f(x)p_n(x) \dx} &\le
    \int_{a}^{b} |f(x)| |f(x) - p_n(x)| \dx \\ &\le
    \frac{\abs{b-a}}{n} \int_{a}^{b} |f(x)| \dx.
  \end{align*}
  Since the last integral is constant and this is true for any $n \geq 1$
  we get that $\int_{a}^{b} f(x)f(x) \dx$ as wanted.

  We can now go back to the original exercise.
  In order to show that $f$ is an even function we will
  show that $f(x) - f(-x) = 0$.
  As we have shown earlier it suffices to show that
  $\int_{-1}^{1} \del{f(x) - f(-x)} x^n \dx = 0$ for every $n \geq 0$.
  From what it given, for every odd positive $n$ we have
  \[
    \int_{-1}^{1} \del{f(x) - f(-x)} x^n \dx =
    \underbrace{\int_{-1}^{1} f(x)  x^n \dx}_{=\,0} -
    \underbrace{\int_{-1}^{1} f(-x) x^n \dx}_{=\,0} = 0.
  \]
  For any even $n \geq 0$ we have that $x^n$ is an even function and thus
  \[
    \int_{-1}^{1}
    \underbrace{\del{f(x) - f(-x)} x^n}_{\text{odd function}} \dx = 0
  \]
  which completes the proof.
\end{solution}

\begin{exercise}
  Consider $C_{\R}([a,b])$, with the metric $d(f,g) = \norm{f - g}_{\infty}$.
  Use the Weierstrass approximation theorem to show
  that $C_{\R}([a,b])$ is separable.
\end{exercise}
\begin{solution}
  In this solution we consider $\Q[x]$ and $\R[x]$ as the rings of polynomials
  with rational and real coefficients respectively but restricted
  to the domain $[a,b]$.
  We know from set theory that $\Q[x]$ is a countable set.
  We will show that $\Q[x]$ is dense in $\C_{\R}([a,b])$.
  First, we notice that $\Q[x]$ is dense in $\R[x]$ since
  given $r(x) = r_0 + r_1 x + \cdots + r_n x^n \in \R[x]$ we know
  that there exists a sequence of rational numbers
  $\set{q_{i,m}}_{m=1}^{\infty}$ that converges to $r_i$ for all
  $0 \le i \le n$.
  Let $q_k(x) = \sum_{i=0}^{n} q_{i,m} x^i$ be a polynomial sequence in $\Q[x]$.
  Denote $g_{i,m} := r_i - q_{i,m}$.
  We now have
  \[
    \abs{r(x) - q_m(x)} =
    \abs{\sum_{i=0}^{n} g_{i,m} x^i} \le
    \sum_{i=0}^{n} \abs{g_{i,m}} \abs{x^i}
    \taking{m \to \infty} 0,
  \]
  where the last step follows from the fact this is a finite sum
  of products of $\abs{g_{i,m}}$ which goes to zero, and $\abs{x^i}$
  which is bounded in $[a,b]$.
  This implies that $\R[x] \subseteq \overline{\Q[x]}$
  which means that
  $\overline{\R[x]} \subseteq \overline{\overline{\Q[x]}} = \overline{\Q[x]}$.
  The Weierstrass approximation theorem gives us that
  $C_{\R}([a,b]) \subseteq \overline{\R[x]}$.
  It follows that
  $C_{\R}([a,b]) \subseteq \overline{\Q[x]}$
  so $C_{\R}([a,b]) = \overline{\Q[x]}$ which completes the proof.
  % another solution
  % Let us define the set of polygonal rational function on $[a,b]$
  % with factor $n$ as
  % the set of all continuous functions $f$ on $[a,b]$ of the form
  % $f(x) = q_n$ where $x = x_i$ for some partition $\set{x_i}_{i=1}^{n}$
  % of $[a,b]$, and so that $f$ is linear on $[x_i,x_{i+1}]$ for
  % $1 \le i < n$.
  % We denote this set $Q_n$.
  % Let $f \in C_{\R}([a,b])$.
  % Then since $[a,b]$ is compact $f$ is uniformly continuous so there
  % exists a sufficienly large $n$ so that $|f(x) - f(y)| < \epsilon$
  % whenever $|x - y| < 1/n$.
  % We can now define a function $g$ by $g(k/n) = f(k/n)$ for $k = 0,1,\dots,n$,
  % and $g$ linear on each interval $(k/n,(k+1)/n)$.
  % From the definition of uniform continuity
  % we necessarily have that $\norm{f - g}_{\infty} \le \epsilon$.
  % We now choose $h \in Q_n$ so that $|h(k/n) - g(k/n)| < \epsilon$
  % for each $k$.
  % It follows that $\norm{g - h}_{\infty} < \epsilon$
  % and thus $\norm{f - h}_{\infty} < 2 \epsilon$.
  % This means that $Q_n$ is dense in $C_{\R}([a,b])$ and we are done.
\end{solution}

\begin{exercise}
  Let $X$ be a compact metric space.
  Prove that $C_{\R}(X)$ is seperable.
  % use the stone-weierstrass theorem along with
  % 1 a compact metric space is seperatble | check
  % 2 a compact metric space is normal | 
  % 3 uryshohn's lemma
\end{exercise}
\begin{solution}
  We know that when $X$ is a compact metric space, it is seperable.
  Let $\set{x_n}_{n=1}^{\infty}$ be a countable set of points dense in $X$.
  We also know that a compact metric space $X$ is normal,
  and from Urysohn's lemma we get that any two disjoint closed subsets can be
  separated by a continuous function.
  This means that any for every closed and disjoint $A,B \subset X$
  there exists a continuous function $f \colon X \to [0,1]$
  from $X$ into the unit interval $[0,1]$ such that $f(a) = 0$
  for all $a \in A$ and $f(b) = 1$ for all $b \in B$.
  Since $X$ is a compact metric space it has a countable open base
  $\mathcal B := \set{U_n}_{n \in \N}$.
  We will choose the standard base with open balls of raional radii.
  For each pair $(n,m)$ such that $\overline{U_n} \subseteq U_m$
  we can generate a function (by Urysohn's lemma) $f_{n,m}$
  so that $f(\overline{U_n}) = \set{0}$ and
  $f\del{X \setminus U_m} = \set{1}$.
  Let $\mathcal A \subseteq C(X)$ be the subalgebra generated
  by the functions $\set{f_{n,m}}_{n,m \in \N}$ and the constant
  function $g$ given by $g(X) = \set{1}$, under $\R$.
  Since $c \cdot g \in \mathcal A$ for any $c \in \Q$, the constant functions
  are contained in the closure of this subalgebra.
  Let $x$, $y$ be different points in $X$.
  Since $X$ is a metric space, there exist open balls
  $B_n = B(x,\delta)$, $B_m = B(x,2\delta)$ so that $\delta \in \Q$.
  We have that $B_n,B_m \in \mathcal B$ and that $\overline{B_n} \subseteq B_m$.
  This implies that there exists $f_{n,m} \in \mathcal A$ so that
  $f(x) = 0$ and $f(1)$ which implies that $\mathcal A$ seperates points.
  From the Stone--Weierstrass theorem $\overline{\mathcal A}$
  is dense in $C_{\R}(X)$,
  which implies that $\mathcal A$ is dense in $C_{\R}(X)$,
  but again, $\mathcal A$ is countable so $C_{\R}(X)$ is seperable which
  completes the proof.
\end{solution}

\newpage

\section{Hilbert spaces}

\begin{exercise}
  Let $H$ be a Hilbert space.
  \begin{enumerate}
    \item[(1)] Let $A \subset H$ be an arbitrary subset.
      Prove that $A^{\perp}$ is a closed subspace of $H$.
    \item[(2)] Let $A \subset H$ be an arbitrary subset.
      Find $A \cap A^{\perp}$.
    \item[(3)] Let $A \subset H$ be an arbitrary subset.
      Prove that $A^{\perp} = \overline A^{\perp}$.
  \end{enumerate}
\end{exercise}
\begin{solution} \phantom{}
  \begin{enumerate}
    \item[(1)] By definition
      \[
        A^{\perp} = \set{h \in H \mid \ip{h}{a} = 0 \text{ for all } a \in A}.
      \]
      Let $\set{x_n}_{n=1}^{\infty}$ be a sequence in $A^{\perp}$
      that converges to $x \in H$.
      Since $\ip{\cdot}{\cdot}$ is a continuous function we have
      for all $a \in A$ that
      \[
        0 =
        \ip{0}{a} =
        \lim_{n \to \infty} \ip{x - x_n}{a} =
        \lim_{n \to \infty} \ip{x}{a} - \ip{x_n}{a} =
        \ip{x}{a},
      \]
      which implies $x \in A^{\perp}$.
    \item[(2)] We obtain from the definition that
      \[
        A \cap A^{\perp} =
        \set{h \in H \cap A \mid \ip{h}{a} = 0 \text{ for all } a \in A} =
        \set{h \in A \mid \ip{h}{a} = 0 \text{ for all } a \in A}.
      \]
      We notice that in particular
      \[
        \set{h \in A \mid \ip{h}{a} = 0 \text{ for all } a \in A} \subset
        \set{h \in A \mid \ip{h}{h} = 0} \subseteq \set{0}.
      \]
      This implies that
      \[
        A \cap A^{\perp} =
        \begin{cases}
          \set{0}, &0 \in A, \\
          \emptyset, &0 \notin A.
        \end{cases}
      \]
    \item[(3)] From the definition and since $A \subset \overline A$ we have
      \[
        \overline A^{\perp} =
        \set{h \in H \mid \ip{h}{a} = 0 \text{ for all } a \in \overline A}
        \subseteq
        \set{h \in H \mid \ip{h}{a} = 0 \text{ for all } a \in A} = A^{\perp},
      \]
      so it suffices to show $A^{\perp} \subseteq \overline A^{\perp}$
      in order to complete the proof.
      Let $a \in A^{\perp}$.
      For all $y \in \overline A$ there exists a sequence
      $\set{x_n}_{n=1}^{\infty}$ in $A$ that converges to $y$.
      The inner product is continuous so we get,
      \[
        \ip{a}{y} =
        \lim_{n \to \infty} \ip{a}{x_n} =
        \lim_{n \to \infty} 0 = 0,
      \]
      which implies $a \in \overline A^{\perp}$ and completes the proof.
  \end{enumerate}
\end{solution}

\begin{exercise}
  Let $H$ be a Hilbert space.
  Let $\set{e_i}_{i=1}^{n}$ be an orthonormal system in $H$.
  Let $M = \vspan\set{e_i}_{i=1}^{n}$.
  Let $P_M$ denote the orthogonal projection of $H$
  onto the closed subspace $M$.
  Prove that for all $h \in H$,
  one has $P_M h = \sum_{i=1}^{n} (h,e_i) e_i$.
\end{exercise}
\begin{solution}
  Denote $m = P_M h$.
  Since $m \in M$ we have that
  \[
    m = \sum_{i=1}^{N} c_i e_i
  \]
  for some constants $c_i$.
  We also notice that
  \[
    \ip{m}{e_k} = \sum_{i=1}^{N} c_i \ip{e_i}{e_k} = c_k,
  \]
  which implies that $m = \sum_{i=1}^{N} \ip{m}{e_k} e_i$.
  We know by a corollary from linear algebra that $h - m$ is orthogonal
  to $M$, which means that $\ip{h - m}{e_i} = 0$ for all $1 \le i \le N$.
  This implies $\ip{h}{e_i} = \ip{m}{e_i}$ which means we can write
  \[
    P_M h := m = \sum_{i=1}^{N} \ip{h}{e_i} e_i
  \]
  which completes the proof.
\end{solution}

\begin{exercise}
  Consider the inner product space $C([0,1])$ of the continuous functions
  on $[0,1]$ with the $\mathcal L^2$ norm:
  $\norm{f}_2 := \sqrt{\int_{0}^{1} |f(x)|^2 \dx}$.
  For $n \geq 2$, let
  \[
    f_n(x) :=
    \begin{cases}
      0, &0 \le x \le \half, \\
      n \del{x - \half}, &\half \le x \le \half + \frac{1}{n}, \\
      1, &\half + \frac{1}{n} \le x \le 1.
    \end{cases}
  \]
  \begin{enumerate}
    \item[(1)] Prove that $\set{f_n}_{n=2}^{\infty}$ is a Cauchy sequence.
    \item[(2)] Prove that $\set{f_n}_{n=2}^{\infty}$ does not converge.
  \end{enumerate}
  % there's a hint
\end{exercise}
\begin{solution} \phantom{}
  \begin{enumerate}
    \item[(1)] Let $m,n \in \N$ so that $m < n$.
      It follows that
      \[
        \norm{f_m - f_n}_2 =
        \sqrt{\int_{0}^{1} |f_m(x) - f_n(x)|^2 \dx} \le
        \sqrt{\int_{0}^{1} \mathbbm{1}_{[1/2,1/2+1/m]} \dx} =
        \frac{1}{\sqrt{m}}.
      \]
      It follows that for $\epsilon > 0$ we can choose $N$ sufficiently
      large so that $\frac{1}{\sqrt{N}} < \epsilon$,
      then for all $n > N$ and $p \in \N$ we have
      \[
        \norm{f_n - f_{n+p}}_2 <
        \frac{1}{\sqrt{n}} <
        \frac{1}{\sqrt{N}} <
        \epsilon,
      \]
      which implies (this is equivalent to the stadard Cauchy criterion)
      that $\set{f_n}_{n=2}^{\infty}$ is a Cauchy sequence.
      
    \item[(2)] Assume, by way of contradiction,
      that $\set{f_n}_{n=2}^{\infty}$
      converges.
      This means that for any
      $\sqrt{\epsilon} > 0$ there exists $N \geq 1$ so that
      for all $n > N$ we have
      \[
        \norm{f - f_n}_2^2 =
        \int_{0}^{1} |f(x) - f_n(x)|^2 \dx < \epsilon.
      \]
      Denote $g_n(x) := f(x) - f_n(x)$.
      It follows that $|g_n(x)|$ converges uniformly to $0$,
      so $g_n(x)$ converges uniformly to $0$,
      so in particular $f_n$ converges pointwise to $f$, and so
      for $y \in (0,1/2)$ and $z \in (1/2,1)$ we get
      \begin{align*}
        f(y) = \lim_{n \to \infty} f_n(y) = 0, \\
        f(z) = \lim_{n \to \infty} f_n(z) = 1,
      \end{align*}
      which implies that
      \begin{align*}
        f(x)\vert_{(0,1/2)} = 0, \\
        f(x)\vert_{(1/2,1)} = 1,
      \end{align*}
      which is a contradiction since $f$ is continuous.
      This contradiction shows that $\set{f_n}_{n=2}^{\infty}$
      does not converge which completes the proof.
  \end{enumerate}
\end{solution}

\begin{exercise}
  Consider the space $PC([0,1])$ of piecewise continuous functions on $[0,1]$
  with the $\mathcal L^2$ norm:
  $\norm{f}_2 := \sqrt{\int_{0}^{1} |f(x)|^2 \dx}$.
  For $n \geq 1$, let
  \[
    f_n(x) =
    \begin{cases}
      0, &0 \le x \le \frac{1}{n}, \\
      x^{-\frac{1}{4}}, &\frac{1}{n} \le x \le 1.
    \end{cases}
  \]
  \begin{enumerate}
    \item[(1)] Prove that $\set{f_n}_{n=1}^{\infty}$ is a Cauchy sequence.
    \item[(2)] Prove that $\set{f_n}_{n=1}^{\infty}$ does not converge.
  \end{enumerate}
  % there's a hint
\end{exercise}
\begin{solution} \phantom{}
  \begin{enumerate}
    \item[(1)] We notice that for $m < n$ we have
      \begin{align*}
        \norm{f_n - f_m}_2 &=
        \sqrt{\int_{0}^{1} |f_n(x) - f_m(x)|^2 \dx} \\ &=
        \sqrt{\int_{0}^{1} |x^{-1/4} \mathbbm{1}_{[1/n,1/m]}|^2 \dx} \\ &=
        \sqrt{\int_{0}^{1} x^{-1/2} \mathbbm{1}_{[1/n,1/m]} \dx} \\ &=
        \sqrt{2\sqrt{x} \Big\vert^{1/m}_{1/n}} \\ &=
        \sqrt{2 \left(\sqrt{1/m} - \sqrt{1/n}\right)} \\ &<
        \sqrt{2 \left(\sqrt{1/m}\right)} \\ &=
        2^{1/2} m^{-1/4}.
      \end{align*}
      For any $\epsilon > 0$ we can choose $N \geq 1$ sufficiently large
      so that $2^{1/2} m^{-1/4} < \epsilon$.
      For all $n,m > N$ so that $m < n$ we have that
      \[
        \norm{f_n - f_m}_2 < 
        2^{1/2} m^{-1/4} <
        2^{1/2} N^{-1/4} < \epsilon
      \]
      which implies that $\set{f_n}_{n=1}^{\infty}$ is a Cauchy sequence.
    \item[(2)] Assume, by way of contradiction, that $g \in PC([0,1])$
      is the limit of $\set{f_n}_{n=1}^{\infty}$.
      Then since $g \in PC([0,1])$ it is bounded.
      Let $M > 0$ be so that $|g(x)| < M$ for $x \in [0,1]$.
      Notice we have that $x^{-1/4} \geq 2M$ for $x \in (0,(2M)^{-4})$,
      which means that
      \[
        0 \le g(x) < M \le \half f_n(x)
      \]
      for $x \in (1/n,(2M)^{-4})$ and $n > 16M^4$.
      This means that for $n > 16M^4$, we have that
      \[
        \norm{f_n - g}_2^2 \geq
        \int_{\frac{1}{n}}^{\frac{1}{16M^4}} |f_n(x) - g(x)|^2 \dx.
      \]
      Using this with the previous inequality we get
      \[
        \norm{f_n - g}_2^2 \geq
        \int_{\frac{1}{n}}^{\frac{1}{16M^4}} |f_n(x) - g(x)|^2 \dx \geq
        \int_{\frac{1}{n}}^{\frac{1}{16M^4}} M^2 \dx > 0,
      \]
      so
      \[
        \norm{f_n - g}_2 \geq
        \sqrt{\int_{\frac{1}{n}}^{\frac{1}{16M^4}} M^2 \dx} > 0,
      \]
      which means that $\lim_{n \to \infty} \norm{f_n - g}_2 > 0$ and thus
      $g$ is not the limit of $\set{f_n}_{n=1}^{\infty}$, which completes the
      proof.
  \end{enumerate}
\end{solution}

\newpage

\section{Fourier series}

\begin{exercise}
  Use the orthonormal basis $\set{e^{2\pi i n x}}_{n=-\infty}^{\infty}$
  on $\mathcal L^2([0,1])$ (with the standard inner product) and apply Bessel's
  inequality to $e^x$ to prove the following identity:
  \[
    \frac{3 - e}{4(e - 1)} =
    \sum_{n=1}^{\infty} \frac{1}{1 + 4 \pi^2 n^2}.
  \]
\end{exercise}
\begin{solution}
  First we compute the Fourier coefficients of $e^x$ which are given by
  \[
    \hat f(n) =
    \int_{0}^{1} e^x e^{-2\pi i n x} \dx =
    \int_{0}^{1} e^{(1 - 2\pi i n) x} \dx =
    \frac{e^{(1 - 2\pi i n) x}}{1 - 2\pi i n} \biggr\vert^{1}_{0} =
    \frac{e^{1 - 2\pi i n} - 1}{1 - 2\pi i n}.
  \]
  In order to use Bessel's inequality we compute $|\hat f(n)|^2$ given by
  \[
    |\hat f(n)|^2 =
    \abs{\frac{e^{1 - 2\pi i n} - 1}{1 - 2\pi i n}}^2 =
    \frac{\abs{e^{1 - 2\pi i n} - 1}^2}{\abs{1 - 2\pi i n}^2}.
  \]
  We can simplify the expression by noticing that $e^{1 - 2\pi i n} = e$
  and that the angle between $1$ and $2 \pi i n$ is $\pi/2$ so
  by using the pythagorian theorem we get
  \[
    |\hat f(n)|^2 =
    \frac{(e - 1)^2}{\abs{1 - 2\pi i n}^2} =
    \frac{(e - 1)^2}{1 + 4\pi^2 n^2}.
  \]
  From the $\mathcal L^2$ convergence of Fourier series we get that
  \[
    \sum_{n \in \Z} |\hat f(n)|^2 =
    \norm{f}_2^2 =
    \int_{0}^{1} e^{2x} \dx =
    \frac{e^{2x}}{2} \biggr\vert^{1}_{0} =
    \frac{e^2 - 1}{2}.
  \]
  From the previous computation we get that
  \[
    \sum_{n \in \Z} \frac{(e - 1)^2}{1 + 4\pi^2 n^2} =
    \frac{e^2 - 1}{2}.
  \]
  After some basic algebraic manipulations
  \[
    (e - 1)^2 \del{1 + 2 \sum_{n=1}^{\infty} \frac{1}{1 + 4\pi^2 n^2}} =
    \frac{e^2 - 1}{2}.
  \]
  Denote
  \[
    x := \sum_{n=1}^{\infty} \frac{1}{1 + 4\pi^2 n^2}
  \]
  We need to solve
  \begin{gather*}
    (e - 1)^2 (1 + 2x) = \frac{e^2 - 1}{2} \\
    \viff \\
    2 (e - 1)^2 (1 + 2x) = e^2 - 1 \\
    \viff \\
    2x = \frac{e^2 - 1 - 2(e - 1)^2}{2 (e - 1)^2}  \\
    \viff \\
    x =
    \frac{4 e - e^2 - 3}{4 (e - 1)^2} =
    \frac{- (e - 1) (e - 3)}{4 (e - 1)^2} =
    \frac{3 - e}{4 (e - 1)}.
  \end{gather*}
  Therefore
  \[
    x := \sum_{n=1}^{\infty} \frac{1}{1 + 4\pi^2 n^2} =
    \frac{3 - e}{4 (e - 1)}.
  \]
\end{solution}

\begin{exercise} \phantom{}
  \begin{enumerate}
    \item[(1)] Find $c \in \C$ for which
      $\inf_{c \in \C} \int_{0}^{1} \abs{e^x - c e^{2\pi i n x}}^2 \dx$
      is attained.
    \item[(2)] Find
      $\inf_{c \in \C} \int_{0}^{1} \abs{e^x - c e^{2\pi i n x}}^2 \dx$.
    \item[(3)] For $n \in \Z$ let
      $H(n) = \inf_{c \in \C} \int_{0}^{1} \abs{e^x - c e^{2\pi i n x}}^2 \dx$.
      Show that $\inf_{n \in \Z} H(n)$ is attained at $n = 0$.
    \item[(4)] Find $\inf_{n \in \Z} H(n)$.
  \end{enumerate}
\end{exercise}
\begin{solution} \phantom{}
  \begin{enumerate}
    \item[(1)] We know that the best approximation for $e^x$ by a complex
      multiple of $e^{2\pi i n x}$ is the projection of $e^x$
      onto $e^{2\pi i n x}$, which is
      \[
        \ip{e^x}{e^{2\pi i n x}} e^{2\pi i n x}.
      \]
      Therefore
      $\inf_{c \in \C} \int_{0}^{1} \abs{e^x - c e^{2\pi i n x}}^2 \dx$
      is attained at
      \[
        c =
        \ip{e^x}{e^{2\pi i n x}} =
        \int_{0}^{1} e^x \overline{e^{2\pi i n x}} \dx =
        \int_{0}^{1} e^x e^{- 2\pi i n x} \dx =
        \hat f(n),
      \]
      for $f(x) = e^x$.
    \item[(2)] We need to find
      \[
        \int_{0}^{1} |e^x - \hat f(n) e^{2\pi i n x}|^2 \dx.
      \]
      Recall $f(x) = e^x$ and denote $g(x) = \hat f(n) e^{2\pi i n x}$.
      With the standard inner product we have that
      \[
        |e^x - \hat f(n) e^{2\pi i n x}|^2 =
        \ip{f - g}{f - g} =
        \ip{f}{f} - \ip{f}{g} - \ip{g}{f} + \ip{g}{g}.
      \]
      For $e_n = e^{2 \pi i n x}$ we have
      \[
        \ip{f}{g} =
        \ip{f}{\hat f(n) e_n} =
        \hat f(n) \ip{f}{e_n} =
        \hat f(n) \overline{\hat f(n)}.
      \]
      where the last equality follows from the definition of Fourier
      coefficients.
      Similarly we get the same result for $\ip{g}{f}$ so
      \[
        \ip{f}{g} = \ip{g}{f} = |\hat f(n)|^2
      \]
      and thus from $\ip{g}{g} = |\hat f(n)|^2$ we get
      \[
        |e^x - \hat f(n) e^{2\pi i n x}|^2 =
        |e^x|^2 + \ip{g}{g} =
        |e^x|^2 - |\hat f(n)|^2.
      \]
      We can now use computations from part $(1)$ and get
      \[
        \int_{0}^{1} |e^x - \hat f(n) e^{2\pi i n x}|^2 \dx =
        \int_{0}^{1} |e^x|^2 - |\hat f(n)|^2 \dx =
        \frac{e^2 - 1}{2} - \frac{(e - 1)^2}{1 + 4 \pi^2 n^2}.
      \]
    \item[(3)] From part $(2)$ we get that
      \[
        H(n) =
        \frac{e^2 - 1}{2} - \frac{(e - 1)^2}{1 + 4 \pi^2 n^2}.
      \]
      Therefore, in order to minimize $H(n)$ we need to minimize
      $1 + 4 \pi^2 n^2$, so clearly since that minimum of
      $1 + 4 \pi^2 n^2$ is attained at $n = 0$ we get that
      $\inf_{n \in \Z} H(n)$ is attained at $n = 0$.
    \item[(4)] By $(3)$ we simply need to find $H(0)$ which is just
      \[
        H(0) =
        \frac{e^2 - 1}{2} - \frac{(e - 1)^2}{1 + 4 \pi^2 0^2} =
        \frac{e^2 - 1}{2} - (e - 1)^2.
      \]
  \end{enumerate}
\end{solution}

\begin{exercise} 
  Solve the following.
  \begin{enumerate}
    \item[(1)] We know that $\{e^{i n x} / \sqrt{2\pi}\}_{n \in \Z}$
      is an orthonormal basis for $\mathcal L^2([-\pi,\pi])$ with the inner product
      \[
        \ip{f}{g} =
        \int_{-\pi}^{\pi} f(x) \overline{g(x)} \dx.
      \]
      So for any real-valued $f \in \mathcal L^2([-\pi,\pi])$ we have
      \[
        f \overset{\mathcal L^2}{=} \sum \hat f(n) \frac{e^{i n x}}{\sqrt{2\pi}},
      \]
      where $\hat f(n) = \langle f,e^{i n x} / \sqrt{2\pi}\rangle$.
      We can then write
      \[
        \hat f(n) =
        \ip{f}{\frac{\cos(n x)}{\sqrt{2\pi}}} +
        i \ip{f}{\frac{\sin(n x)}{\sqrt{2\pi}}}.
      \]
      Show that when plugging this equation into the original $\mathcal L^2$
      equality, we get
      \[
        f \overset{\mathcal L^2}{=}
        \frac{a_0}{2} +
        \sum_{n=1}^{\infty} a_n \cos(nx) + \sum_{n=1}^{\infty} b_n \sin(nx),
      \]
      where
      \begin{align*}
        a_n &= \frac{1}{\pi} \int_{-\pi}^{\pi} f(x) \cos(nx) \dx \\
        b_n &= \frac{1}{\pi} \int_{-\pi}^{\pi} f(x) \sin(nx) \dx
      \end{align*}
      and make sure you understand why $a_0/2$ appears.
      Recall how it was defined.
    \item[(2)] Show that the sequence
      \[
        \set{\frac{1}{\sqrt{2}}, \cos(x), \sin(x), \cos(2x), \sin(2x),\dots}
      \]
      is an orthonormal basis for $\mathcal L^2([-\pi,\pi])$ with the inner product
      \[
        \ip{f}{g} =
        \frac{1}{\pi} \int_{-\pi}^{\pi} f(x) \overline{g(x)} \dx.
      \]
    \item[(3)] Show that if $g$ is an even function,
      then
      \[
        g \overset{\mathcal L^2}{=} a_0/2 + \sum_{n=1}^{\infty} a_n \cos(nx),
      \]
      and if $g$ is an odd function, then
      \[
        g \overset{\mathcal L^2}{=} \sum_{n=1}^{\infty} b_n \sin(nx).
      \]
    \item[(4)] Let $f \in \mathcal L^2([0,\pi])$.
      Extend $f$ to $f_e(x)$ so that $f_e$ is an even function
      in $\mathcal L^2([-\pi,\pi])$.
      Also extend $f$ to $f_o(x)$ so that $f_o$ is an odd function
      in $\mathcal L^2([-\pi,\pi])$.
      Use part $(3)$ and the above extensions to show that $f$ can be
      represented in the two following ways
      \[
        f(x) \overset{\mathcal L^2}{=} \frac{a_0}{2} + \sum_{n=1}^{\infty} a_n \cos(nx),
      \]
      where
      \[
        a_n = \frac{2}{\pi} \int_{0}^{\pi} f(x) \cos(nx) \dx
      \]
      and
      \[
        f(x) \overset{\mathcal L^2}{=} \sum_{n=1}^{\infty} b_n \sin(nx),
      \]
      where
      \[
        b_n = \frac{2}{\pi} \int_{0}^{\pi} f(x) \sin(nx) \dx.
      \]
    \item[(5)] Prove that on $\mathcal L^2([0,\pi])$ with the inner product
      \[
        \ip{f}{g} =
        \frac{2}{\pi} \int_{0}^{\pi} f(x) g(x) \dx,
      \]
      each of the following sequences of functions is an orthonormal basis.
      \begin{enumerate}
        \item[(a)] $\set{1/\sqrt{2},\cos(x),\cos(2x),\cos(3x),\dots}$.
        \item[(b)] $\set{\sin(x),\sin(2x),\sin(3x),\dots}$.
      \end{enumerate}
      Note that the span of $(a)$ and $(b)$ are both not algebras.
      Thus, we can't use the Stone--Weierstrass theorem to prove that each
      of the above sequences is a basis.
      Note that the equation
      \[
        f(x) \overset{\mathcal L^2}{=} \frac{a_0}{2} + \sum_{n=1}^{\infty} a_n \cos(nx),
      \]
      is called the \textit{Fourier cosine series} and equation
      \[
        f(x) \overset{\mathcal L^2}{=} \sum_{n=1}^{\infty} b_n \sin(nx),
      \]
      is called the \textit{Fourier sine series} of $f$.
  \end{enumerate}
\end{exercise}
\begin{solution} \phantom{}
  \begin{enumerate}
    \item[(1)] We have that
      \begin{align*}
        f &\overset{\mathcal L^2}{=}
        \sum
        \del{\ip{f}{\frac{\cos(n x)}{\sqrt{2\pi}}} +
        i \ip{f}{\frac{\sin(n x)}{\sqrt{2\pi}}}}
        \frac{e^{i n x}}{\sqrt{2\pi}} \\ &=
        \sum_{n \in \Z} \ip{f}{\frac{\cos(n x)}{\sqrt{2\pi}}}
          \frac{\cos(n x)}{\sqrt{2\pi}} +
        \sum_{n \in \Z} \ip{f}{\frac{\sin(n x)}{\sqrt{2\pi}}}
        \frac{\sin(n x)}{\sqrt{2\pi}} \\ &+ i\del{
          \sum_{n \in \Z} \ip{f}{\frac{\cos(n x)}{\sqrt{2\pi}}}
          \frac{\sin(n x)}{\sqrt{2\pi}} -
        \sum_{n \in \Z} \ip{f}{\frac{\sin(n x)}{\sqrt{2\pi}}}
        \frac{\cos(n x)}{\sqrt{2\pi}}
        }.
      \end{align*}
      Since $f$ is real valued we get that
      \[
        \sum_{n \in \Z} \ip{f}{\frac{\cos(n x)}{\sqrt{2\pi}}}
          \frac{\sin(n x)}{\sqrt{2\pi}} -
        \sum_{n \in \Z} \ip{f}{\frac{\sin(n x)}{\sqrt{2\pi}}}
          \frac{\cos(n x)}{\sqrt{2\pi}} = 0.
      \]
      Now since $\cos$ is even we get
      \[
        \sum_{n \in \Z} \ip{f}{\frac{\cos(n x)}{\sqrt{2\pi}}}
          \frac{\cos(n x)}{\sqrt{2\pi}} =
        \ip{f}{\frac{\cos(0)}{2}} \cos(0) +
        2 \sum_{n=1}^{\infty} \ip{f}{\frac{\cos(n x)}{\sqrt{2\pi}}}
          \frac{\cos(n x)}{\sqrt{2\pi}}
      \]
      and since $\sin$ is even we get
      \[
        \sum_{n \in \Z} \ip{f}{\frac{\sin(n x)}{\sqrt{2\pi}}}
          \frac{\sin(n x)}{\sqrt{2\pi}} =
        \ip{f}{\frac{\sin(0)}{2\pi}} \sin(0) + 2\sum_{n=1}^{\infty} \ip{f}{\frac{\sin(n x)}{\sqrt{2\pi}}}
          \frac{\sin(n x)}{\sqrt{2\pi}}.
      \]
      Compute the inner products:
      \begin{align*}
        \ip{f}{\frac{\cos(0)}{2}} \cos(0) &= \half \ip{f}{\frac{1}{\pi}} \\ &=
        \frac{1}{\pi} \int_{-\pi}^{\pi} f(x) \dx, \\
        \ip{f}{\frac{\sin(0)}{2}} \sin(0) &= 0.
      \end{align*}
      So by combining these results we get
      \[
        f \overset{\mathcal L^2}{=}
        \frac{a_0}{2} +
        \sum_{n=1}^{\infty} a_n \cos(nx) + \sum_{n=1}^{\infty} b_n \sin(nx),
      \]
      which completes the first part of the exercise.
    \item[(2)] First we check that all vectors are of norm $1$.
      For the constant function:
      \[
        \norm{\frac{1}{\sqrt{2}}}^2 =
        \frac{1}{\pi} \int_{-\pi}^{\pi} \del{\frac{1}{\sqrt{2}}}^2 \dx =
        \frac{1}{\pi} \cdot \half \cdot 2\pi = 1.
      \]
      For $\cos(nx):$
      \[
        \norm{\cos(nx)}^2 =
        \frac{1}{\pi} \int_{-\pi}^{\pi} \cos^2(nx) \dx =
        \frac{1}{\pi} \cdot \pi = 1
      \]
      since the integral is a known result from the proof of orthonormality
      of the standard basis with the standard inner product.
      Similarly,
      \[
        \norm{\sin(nx)}^2 =
        \frac{1}{\pi} \int_{-\pi}^{\pi} \sin^2(nx) \dx =
        \frac{1}{\pi} \cdot \pi = 1.
      \]
      Orthogonality is given immdiately from the fact
      \[
        \ip{f_1}{f_2} = \ip{f_1}{f_2}_{\std}
      \]
      Since
      $\int_{-\pi}^{\pi} \sin(nx) \dx = \int_{-\pi}^{\pi} \cos(nx) \dx = 0$
      we get that $1/\sqrt{2}$ is orthogonal to all elements of the basis.
      We also know from calculus that for $n \neq m$ that
      \[
        \int_{-\pi}^{\pi} \cos(nx) \sin(mx) \dx =
        \int_{-\pi}^{\pi} \sin(nx) \sin(mx) \dx =
        \int_{-\pi}^{\pi} \cos(nx) \cos(mx) \dx = 0.
      \]
      This shows that any two elements of the sequence are orthogonal.
      We have see from part $(1)$ that any $f \in \mathcal L^2([-\pi,\pi])$
      can be written as
      \[
        f \overset{\mathcal L^2}{=}
        \frac{a_0}{2} +
        \sum_{n=1}^{\infty} a_n \cos(nx) + \sum_{n=1}^{\infty} b_n \sin(nx),
      \]
      where
      \begin{align*}
        a_n &= \frac{1}{\pi} \int_{-\pi}^{\pi} f(x) \cos(nx) \dx =
        \ip{f}{\cos(nx)} \\
        b_n &= \frac{1}{\pi} \int_{-\pi}^{\pi} f(x) \sin(nx) \dx =
        \ip{f}{\sin(nx)}
      \end{align*}
      which shows that the set is also complete and completes the proof.
    \item[(3)] Let $g$ be an even function.
      Then since $\sin(nx)$ is odd we get that $g(x) \sin(x)$ is odd so
      \[
        b_n = \frac{1}{\pi} \int_{-\pi}^{\pi} g(x) \sin(nx) \dx = 0
      \]
      so from part $(1)$ we get
      \[
        g \overset{\mathcal L^2}{=} a_0/2 + \sum_{n=1}^{\infty} a_n \cos(nx).
      \]
      Let $g$ be an odd function, then since $\cos(nx)$ is even we
      get that $g(x) \cos(nx)$ is odd so
      \[
        a_n = \frac{1}{\pi} \int_{-\pi}^{\pi} f(x) \cos(nx) \dx = 0
      \]
      so from part $(1)$ we get
      \[
        g \overset{\mathcal L^2}{=} \sum_{n=1}^{\infty} b_n \sin(nx).
      \]
    \item[(4)] We define the extension $f_e(x)$ of $f(x)$ to be even.
      From $(3)$ we know that if $f_e(x)$ is even then we can write
      \[
        f_e(x) \overset{\mathcal L^2}{=} a_0/2 + \sum_{n=1}^{\infty} a_n \cos(nx).
      \]
      Recall that since $f_e(x)$ and $\cos(nx)$ are even,
      their product $f_e(x) \cos(nx)$ is even.
      Thus by definition of $a_n$ and properties of the integral
      \[
        a_n =
        \frac{1}{\pi} \int_{-\pi}^{\pi} f_e(x) \cos(nx) \dx =
        \frac{2}{\pi} \int_{0}^{\pi} f_e(x) \cos(nx) \dx.
      \]
      We define the extension $f_o(x)$ of $f(x)$ to be odd.
      From $(3)$ we know that if $f_o(x)$ is odd then we can write
      \[
        f_o(x) \overset{\mathcal L^2}{=} \sum_{n=1}^{\infty} b_n \sin(nx).
      \]
      Recall that since $f_o(x)$ and $\sin(nx)$ are even,
      their product $f_o(x) \sin(nx)$ is even.
      Thus by definition of $b_n$ and properties of the integral
      \[
        b_n =
        \frac{1}{\pi} \int_{-\pi}^{\pi} f_o(x) \sin(nx) \dx =
        \frac{2}{\pi} \int_{0}^{\pi} f_o(x) \sin(nx) \dx.
      \]
      This completes this part of the exercise.
    \item[(5)] We will verify orthonormality of completeness seperately,
      but both are immediate results of previous part of the exercise.
      \begin{enumerate}
        \item[(a)] We are only considering functions $f$ in $\mathcal L^2([0,\pi])$,
          so we can consider their even extension $f_e$ in $\mathcal L^2([-\pi,\pi])$.
          In part $(4)$ we saw that we can write
          \[
            f_e(x) \overset{\mathcal L^2}{=} a_0/2 + \sum_{n=1}^{\infty} a_n \cos(nx).
          \]
          for
          \[
            a_n =
            \frac{2}{\pi} \int_{0}^{\pi} f_e(x) \cos(nx) \dx,
          \]
          so simply by restricting all the functions to the domain
          $[0,\pi]$ we get that set $(a)$ is complete.
          We can also notice that $(a)$ is contained in
          \[
            \set{\frac{1}{\sqrt{2}}, \cos(x), \sin(x), \cos(2x), \sin(2x),\dots}
          \]
          which is orthonormal with respect to the inner product in
          part $(2)$ given by
          \[
            \ip{f}{g}_{(2)} =
            \frac{1}{\pi} \int_{-\pi}^{\pi} f(x) \overline{g(x)} \dx.
          \]
          Thus in particular set $(a)$ is orthonormal with respect to
          $\ip{f}{g}_{(2)}$.
          However, for even, real-valued functions $f$, $g$ we see that
          \[
            \ip{f}{g}_{(2)} =
            \frac{1}{\pi} \int_{-\pi}^{\pi} f(x) \overline{g(x)} \dx =
            \frac{2}{\pi} \int_{0}^{\pi} f(x) g(x) \dx =
            \ip{f}{g}.
          \]
          This means that $(a)$ is orthonormal with respect to the inner
          product given in this part, so it is an orthonormal basis.
        \item[(b)] We are only considering functions $f$ in $\mathcal L^2([0,\pi])$,
          so we can consider their odd extension $f_o$ in $\mathcal L^2([-\pi,\pi])$.
          In part $(4)$ we saw that we can write
          \[
            f_o(x) \overset{\mathcal L^2}{=} \sum_{n=1}^{\infty} b_n \sin(nx).
          \]
          for
          \[
            b_n =
            \frac{2}{\pi} \int_{0}^{\pi} f_o(x) \sin(nx) \dx,
          \]
          so simply by restricting all the functions to the domain
          $[0,\pi]$ we get that set $(b)$ is complete.
          We can also notice that $(b)$ is contained in
          \[
            \set{\frac{1}{\sqrt{2}}, \cos(x), \sin(x), \cos(2x), \sin(2x),\dots}
          \]
          which is orthonormal with respect to the inner product in
          part $(2)$ given by
          \[
            \ip{f}{g}_{(2)} =
            \frac{1}{\pi} \int_{-\pi}^{\pi} f(x) \overline{g(x)} \dx.
          \]
          Thus in particular set $(b)$ is orthonormal with respect to
          $\ip{f}{g}_{(2)}$.
          However, for odd, real-valued functions $f$, $g$ we see that
          \[
            \ip{f}{g}_{(2)} =
            \frac{1}{\pi} \int_{-\pi}^{\pi} f(x) \overline{g(x)} \dx =
            \frac{2}{\pi} \int_{0}^{\pi} f(x) g(x) \dx =
            \ip{f}{g}.
          \]
          This means that $(b)$ is orthonormal with respect to the inner
          product given in this part, so it is an orthonormal basis.
      \end{enumerate}
  \end{enumerate}
\end{solution}

\begin{exercise}
 Use Abel's summation formula to prove that if $\set{b_n}_{n=1}^{\infty}$ is
 a sequence so that the sequence $B_n = \sum_{j=1}^{n} b_n$ is bounded,
 and $\set{a_n}_{n=1}^{\infty}$ is monotonically nonincreasing with
 $\lim_{n \to \infty} a_n = 0$, then the series
 $\sum_{n=1}^{\infty} a_n b_n$ is convergent.
\end{exercise}
\begin{solution}
  Since the sequence of partial sums $B_n$ is bounded there exists $M > 0$
  so that for all $n \geq 0$ we have $|B_n| < M$.
  Consider Abel's summation formula (summation by parts formula):
  \[
    \sum_{i=1}^{n} a_i b_i =
    a_n B_n - \sum_{i=1}^{n-1} B_i (a_{i+1} - a_i).
  \]
  We see that the series $\sum_{i=1}^{\infty} B_i (a_{i+1} - a_i)$
  converges since for all $n$
  \begin{align*}
    \sum_{i=1}^{n} \abs{B_i (a_{i+1} - a_i)} &=
    \sum_{i=1}^{n} \abs{B_i} \cdot \abs{(a_{i+1} - a_i)} \le
    M \sum_{i=1}^{n} \abs{(a_{i+1} - a_i)} \\ &=
    M \sum_{i=1}^{n} (a_{i} - a_{i+1}) = M (a_1 - a_{n + 1}).
  \end{align*}
  Since $\lim_{n \to \infty} a_n = 0$ we get that
  $M (a_1 - a_{n + 1})$ is bounded,
  so $\sum_{i=1}^{n} \abs{B_i (a_{i+1} - a_i)}$ is bounded,
  but since $\sum_{i=1}^{n} \abs{B_i (a_{i+1} - a_i)}$
  is a bounded series of nonnegative numbers, it must converge.
  Therefore $\sum_{i=1}^{\infty} B_i (a_{i+1} - a_i)$ converges
  absolutely, and in particular it converges.
  The sequence $a_n B_n$ converges to $0$ since $\lim_{n \to \infty} a_n = 0$
  and $B_n$ is bounded.
  Therefore
  \[
    \lim_{n \to \infty} a_n B_n - \sum_{i=1}^{n-1} B_i (a_{i+1} - a_i)
  \]
  exists so
  \[
    \lim_{n \to \infty} \sum_{i=1}^{n} a_i b_i =
    \lim_{n \to \infty} a_n B_n - \sum_{i=1}^{n-1} B_i (a_{i+1} - a_i)
  \]
  exists which completes the proof.
\end{solution}

\begin{exercise}
  Let $f \in \mathcal L^2([-\pi,\pi])$ be given by
  \[
    f(x) = \mathbbm{1}_{[a,b]}(x) =
    \begin{cases}
      1, & x \in [a,b], \\
      0, & x \in [-\pi,\pi] \setminus [a,b],
    \end{cases}
  \]
  for $-\pi < a < b < \pi$.
  \begin{enumerate}
    \item[(1)] Show that the Fourier series of $f$ is given by
      \[
        f \overset{\mathcal L^2}{=}
        \frac{b - a}{2\pi} +
        \sum_{n \neq 0} \frac{e^{- i n a} - e^{- i n b}}{2 \pi i n} e^{i n x},
      \]
      where $\sum_{n \neq 0}$ is summing over $n \in \Z \setminus \set{0}$.
    \item[(2)] Show that the Fourier series on the right hand side above
      converges at every point $x \in [-\pi,\pi] \setminus \set{a,b}$
      by using Abel's summation.
    \item[(3*)] Show that the Fourier series on the right hand
      side above does not converge when $x = a$ or $x = b$.
  \end{enumerate}
  Note that by Dirichlet’s Theorem, the symmetric partial sums
  $\sum_{n=-N}^{N} \frac{e^{- i n a} - e^{- i n b}}{2 \pi i n} e^{i n x}$
  converge for all $x \in [-\pi,\pi]$ as $N \to \infty$.
\end{exercise}
\begin{solution} \phantom{}
  \begin{enumerate}
    \item[(1)] First we calculate the
      Fourier coefficients $\set{c_n}_{z \in \Z}$.
      For $n = 0$ we get
      \[
        c_0 =
        \ip{f}{\frac{e^{i 0 x}}{\sqrt{2\pi}}} =
        \ip{f}{\frac{1}{\sqrt{2\pi}}} =
        \frac{1}{\sqrt{2\pi}} \int_{-\pi}^{\pi} \mathbbm{1}_{[a,b]}(x) \dx =
        \frac{b - a}{\sqrt{2\pi}}.
      \]
      For $n \neq 0$ we get
      \[
        c_n =
        \ip{f}{\frac{e^{i n x}}{\sqrt{2\pi}}} =
        \frac{1}{\sqrt{2\pi}}
          \int_{-\pi}^{\pi} \mathbbm{1}_{[a,b]}(x) e^{- i n x} \dx =
        \frac{1}{\sqrt{2\pi}} \int_{a}^{b} e^{- i n x} \dx.
      \]
      Therefore
      \[
        c_n =
        \frac{1}{\sqrt{2\pi}} \int_{a}^{b} e^{- i n x} \dx =
        \frac{1}{\sqrt{2\pi}} \eval{\frac{e^{-i n x}}{-i n}}_{a}^{b} =
        \frac{e^{-i n b} - e^{-i n a}}{-i n \sqrt{2\pi}} =
        \frac{e^{-i n a} - e^{-i n b}}{i n \sqrt{2\pi}}.
      \]
      Now by the Fourier series we get
      \[
        f \overset{\mathcal L^2}{=}
        \frac{b - a}{\sqrt{2\pi}} \cdot \frac{1}{\sqrt{2\pi}} +
        \sum_{n \neq 0} \frac{e^{-i n a} - e^{-i n b}}{\sqrt{2\pi} i n} \cdot
        \frac{e^{i n x}}{\sqrt{2\pi}} =
        \frac{b - a}{2\pi} +
        \sum_{n \neq 0} \frac{e^{-i n a} - e^{-i n b}}{2\pi i n}
        e^{i n x}
      \]
      which completes the proof.
    \item[(2)] (bad solution) Let $x \in [-\pi,\pi] \setminus \set{a,b}$. %to be added %to fix %problem
      Write the partial sums as
      \[
        S_N =
        \frac{b - a}{2 \pi} +
        \sum_{0 < |n| \le N} \frac{e^{-i n a} - e^{-i n b}}{2\pi i n} e^{i n x}=
        \frac{b - a}{2 \pi} +
        \sum_{0 < |n| \le N} \frac{e^{i n (x - a)}}{2\pi i n} -
        \sum_{0 < |n| \le N} \frac{e^{i n (x - b)}}{2\pi i n}.
      \]
      In order to show that this sequence converges it suffices to show that
      \[
        \sum_{0 < |n| \le N} \frac{e^{i n (x - a)}}{2\pi i n} \tand
        \sum_{0 < |n| \le N} \frac{e^{i n (x - b)}}{2\pi i n}
      \]
      converge.
      However these sums are by definition equal to the Dirichlet's kernel
      and we know that
      \[
        \sum_{0 < |n| \le N} e^{i n \theta} =
        D_N(\theta) =
        \frac{\sin\del{(N + \half) \theta}}{\sin(\theta/2)}.
      \]
      Therefore
      \[
        \abs{\sum_{0 < |n| \le N} e^{i n (x - a)}} =
        \abs{D_N(x-a)} \le
        \frac{1}{\abs{\sin\del{\frac{x - a}{2}}}}.
      \]
      This bounded is indepent of $n$ and well defined when $x \neq a$.
      Now since we can take the constant $1/2\pi i$ out of the sum
      we get that
      \[
        \sum_{0 < |n| \le N} \frac{e^{i n (x - a)}}{2\pi i n} =
        \frac{1}{2\pi i}\sum_{0 < |n| \le N} \frac{e^{i n (x - a)}}{n}
      \]
      converges from Dirichlet's test
      because $\sum_{0 < |n| \le N} e^{i n (x - a)}$ is bounded and
      $1/n$ is decreasing.
      There, after repeating the same argument for the second sum
      we get that $S_N$ converges for $x \notin \set{a,b}$
      which completes the proof.
    \item[(3)] (bad solution) Let $x = a$. %to be added
      Then we have
      \[
        S_N =
        \frac{b - a}{2 \pi} +
        \sum_{0 < |n| \le N} \frac{1 - e^{- i n (b - a)}}{2\pi i n}.
      \]
      We will show that the latter term does not converge.
      \[
        \sum_{0 < |n| \le N} \frac{1 - e^{- i n (b - a)}}{2\pi i n} =
        \sum_{0 < |n| \le N} \frac{1}{2\pi i n} -
        \sum_{0 < |n| \le N} \frac{e^{- i n (b - a)}}{2\pi i n} =
        - \sum_{0 < |n| \le N} \frac{e^{- i n (b - a)}}{2\pi i n}
      \]
      because $\sum_{0 < |n| \le N} \frac{1}{2\pi i n} = 0$ since
      $1/2 \pi i n$ is an odd function in $n$.
      Notice that when we change $n$ to $-n$ we get
      \[
        - \sum_{0 < |n| \le N} \frac{e^{- i n (b - a)}}{2\pi i n} =
        \sum_{0 < |n| \le N} \frac{e^{i n (b - a)}}{2\pi i n}.
      \]
      We also have
      \[
        \sum_{0 < |n| \le N} \frac{e^{i n (b - a)}}{2\pi i n} =
        \sum_{n = 1}^{N} \del{
        \frac{e^{i n (b - a)}}{i n} + \frac{e^{i n (b - a)}}{- i n}} =
        \sum_{n = 1}^{N} \frac{2\sin\del{n (b - a)}}{n}
      \]
      where the last transition is from the identity
      $e^{i \theta} - e^{-i \theta} = 2i \sin(\theta)$.
      Thus, the original sum becomes
      \[
        S_N =
        \frac{b - a}{2 \pi} +
        \frac{1}{2 \pi}
        \sum_{n = 1}^{N} \frac{2\sin\del{n (b - a)}}{n} =
        \frac{b - a}{2 \pi} +
        \frac{1}{\pi}
        \sum_{n = 1}^{N} \frac{\sin\del{n (b - a)}}{n}.
      \]
      We need to show that
      \[
        \sum_{n = 1}^{N} \frac{\sin\del{n (b - a)}}{n}
      \]
      does not converge.
      Similarly we can show this for $x = b$ and thus this part is complete.
  \end{enumerate}
\end{solution}

\newpage

\section{Weak convergence}

\begin{exercise}
  Prove that the weak limit is unique.
  That is, if $h_n \taking{w} h_1$ and $h_n \taking{w} h_2$,
  then $h_1 = h_2$.
\end{exercise}
\begin{solution}
  Let $h_1,h_2 \in H$ be
  such that $h_n \taking{w} h_1$ and $h_n \taking{w} h_2$.
  Let $g \in H$, then from weak convergence we get that
  \begin{gather*}
    \lim_{n \to \infty} \ip{h_n}{g} = \ip{h_1}{g} \\
    \lim_{n \to \infty} \ip{h_n}{g} = \ip{h_2}{g}
  \end{gather*}
  and from uniqueness of the limit in $\C$ we get that
  \[
    \ip{h_1}{g} = \ip{h_2}{g} \implies
    \ip{h_1 - h_2}{g} = 0.
  \]
  Since the choice of $g$ was arbitrary this implies that
  \[
    \forall g \quad \ip{h_1 - h_2}{g} = 0 \implies h_1 = h_2,
  \]
  which completes the proof.
\end{solution}

\begin{exercise}
  Let $T \in B(H,K)$.
  Then $\norm{T} = \sup_{\substack{h \in H,k \in K \\
  \norm{h} = \norm{k} = 1}} \abs{\ip{T(h)}{k}}$.
\end{exercise}
\begin{solution}
  Let $h \in H$ and $k \in K$ be so that $\norm{h} = \norm{k} = 1$.
  Then we have from the Cauchy--Schwartz inequality that
  \[
    \abs{\ip{T(h)}{k}} \le
    \norm{Th} \cdot \underbrace{\norm{k}}_{=\,1} \le
    \norm{T} \cdot \norm{h} = \norm{T}.
  \]
  On the other hand, since for any $h \in H$ so that $\norm{h} = 1$
  we can choose $k = Th / \norm{Th}$ (which also has norm $1$).
  We get that
  \[
    \sup_{\substack{h \in H,k \in K \\
      \norm{h} = \norm{k} = 1}} \abs{\ip{T(h)}{k}} \geq
    \sup_{\substack{h \in H \\ \norm{h} = 1}}
      \abs{\ip{T(h)}{\frac{T(h)}{\norm{Th}}}} =
    \sup_{\substack{h \in H \\ \norm{h} = 1}}
    \frac{1}{\norm{Th}} \cdot \norm{Th}^2 = \norm{T}. 
  \]
  From these two inequalities we get that
  \[
    \norm{T} = \sup_{\substack{h \in H,k \in K \\
    \norm{h} = \norm{k} = 1}} \abs{\ip{T(h)}{k}},
  \]
  which completes the proof.
\end{solution}

\begin{exercise} \phantom{}
  \begin{enumerate}
    \item[(1)] Let $h_n \taking{w} h$ and assume
      that $l := \lim_{n \to \infty} \norm{h_n}$ exists.
      Show that $\norm{h} \le l$.
    \item[(2)] More generally, show that if $h_n \taking{w} h$,
      then $\norm{h} \le \liminf_{n \to \infty} \norm{h_n}$.
  \end{enumerate}
\end{exercise}
\begin{solution} \phantom{}
  \begin{enumerate}
    \item[(1)] By definition of weak convergence for $h \in H$ we have
      \[
        \lim_{n \to \infty} \ip{h_n}{h} =
        \ip{h}{h} =
        \abs{\ip{h}{h}} =
        \norm{h}^2.
      \]
      We also know that
      \[
        \norm{h_n - h}^2 =
        \norm{h_n}^2 - 2\Re(\ip{h_n}{h}) + \norm{h}^2
      \]
      so we can write
      \[
        \norm{h_n}^2 =
        \norm{h_n - h}^2 + 2\Re(\ip{h_n}{h}) - \norm{h}^2.
      \]
      Note that we already know some of those values when we take $n$
      to infinity:
      \[
        \norm{h_n}^2 \taking{n \to \infty} l^2 \tand
        \text{(from weak convergence) }
        \Re(\ip{h_n}{h}) \taking{n \to \infty} \norm{h}^2.
      \]
      Therefore
      \[
        0 \le
        \norm{h_n}^2 \taking{n \to \infty}
        l^2 - 2\norm{h}^2 + \norm{h^2} = l^2 - \norm{h}^2.
      \]
      This implies that $l^2 \geq \norm{h}^2$ and since both $l$
      and $\norm{h}$ are nonnegative this implies that $\norm{h} \le l$.
    \item[(2)] Denote $l := \liminf_{n \to \infty} \norm{h_n}$.
      When $l = \infty$ the identity clearly holds.
      When $l < \infty$ there exists a subsequence of $h_n$ so that
      \[
        \lim_{k \to \infty} \norm{h_{n_k}} = l.
      \]
      Let $g \in H$.
      Since $\lim_{n \to \infty} \ip{h_n}{g} = \ip{h}{g}$,
      we also have $\lim_{n \to \infty} \ip{h_{n_k}}{g} = \ip{h}{g}$.
      The choice of $g$ was arbitrary so $h_{n_k} \taking{w} h$.
      From $(1)$ we get that $\norm{h} \le l$ which completes the proof.
  \end{enumerate}
\end{solution}

\begin{exercise}
  A not necessarily norm-continuous
  (or equivalently, not necessarily bounded)
  linear operator from a Hilbert space $H$ to a Hilbert space $K$
  is called weakly sequentially continuous if $h_n \taking{w} h$
  implies $Ah_n \taking{w} Ah$.
  In class we will prove that if $A$
  is weakly sequentially continuous,
  then it is norm-continuous (or equivalently, bounded).
  Show that if $A$ is norm-continuous
  (or equivalently, bounded)
  then it is weakly sequentially continuous.
\end{exercise}
\begin{solution}
  Let $A \in B(H,K)$
  and let $h_n \taking{w} h$.
  We need to show that for any $g \in H$ we have
  \[
    \lim_{n \to \infty} \ip{Ah_n}{g}_K = \ip{Ah}{g}_K.
  \]
  Since $A \in B(H,K)$ is has a conjugate $A^*$.
  Thus, by $h_n \taking{w} h$ we get
  \[
    \ip{Ah_n}{g}_K =
    \ip{h_n}{A^* g}_H \taking{n \to \infty}
    \ip{h}{A^* g}_H =
    \ip{Ah}{g}_H,
  \]
  which completes the proof.
\end{solution}

\newpage

\section{Bounded operators}

\begin{exercise}
  Let $A \in B(H,K)$,
  where $H,K$ are Hilbert spaces.
  Show that if the operators $A$ and $A^*$
  are both bounded from below, then $A$ and $A^*$ are invertible.
\end{exercise}
\begin{solution}
  Since $A$ is bounded from below there exists $c > 0$ so that
  for all $0 \neq h \in H$ we have
  \[
    \norm{A h}_{K} \geq c \cdot \underbrace{\norm{h}_{H}}_{> 0} > 0.
  \]
  Therefore $\ker(A) = 0$, which implies that $A$ is injective.
  An identical argument shows that $\ker(A^*) = 0$ is also injective
  (and thus by the Cantor--Schr\"oder--Bernstein $\abs{H} = \abs{K}$).
  Since the kernels are trivials the kernel-image identites in Hilbert
  spaces give
  \begin{gather*}
    \set{0_H}^{\perp} = \ker(A)^{\perp} = \overline{\im(A^*)}, \\
    \set{0_K}^{\perp} = \ker(A^*)^{\perp} = \overline{\im(A)}.
  \end{gather*}
  Since operators bound from below are closed maps we get in particular
  that $\im(A)$ and $\im(A^*)$ are closed and so
  \begin{gather*}
    H = \set{0_H}^{\perp} = \im(A^*), \\
    K = \set{0_K}^{\perp} = \im(A).
  \end{gather*}
  This implies that $A$ and $A^*$ are both surjective.
  All bounded operators which are bijections are invertible,
  which implies that $A$ and $A^*$ both are invertible and completes the proof.
\end{solution}

\begin{exercise}
  \label{ex:lemma-of-w}
  Let $\set{W_n}_{n=1}^{\infty}$ be a sequence of positive numbers
  satisfying $\sup_{n \geq 1} W_n = \infty$.
  In this exercise we will show that
  there exists a sequence $\set{a_n}_{n=1}^{\infty}$ of nonnegative
  numbers such that
  \[
    \sum_{n=1}^{\infty} a_n < \infty \tand
    \sum_{n=1}^{\infty} a_n W_n = \infty.
  \]
  \begin{enumerate}
    \item[(1)] Show that by the assumption on $\set{W_n}_{n=1}^{\infty}$,
      there exists a subsequence $\set{W_{n_k}}_{k=1}^{\infty}$
      such that $W_{n_{k}} \geq k$.
    \item[(2)] Define
      \[
        a_n :=
        \begin{cases}
          \frac{1}{k^2}, &\text{if } n = n_k \text{ for } k=1,2,\dots; \\
          0, &\text{otherwise}.
        \end{cases}
      \]
      Show that $\set{a_n}_{n=1}^{\infty}$ is the desired sequence.
  \end{enumerate}
\end{exercise}
\begin{solution} \phantom{}
  \begin{enumerate}
    \item[(1)] From $\sup_{n \geq 1} W_n = \infty$ we know that there exists
      a subsequence $W_{n_{j}}$
      so that $W_{n_{j}} \taking{j \to \infty} \infty$.
      Then by definition we know that for all $\N \ni k > 0$ there exists
      $j_k \in \N$ so that for all $j > j_k$ we have $W_{n_{j}} > k$.
      This allows us to define a sequence subsequence of $W_{n_j}$
      and in particular of $W_{n}$ which we can denote
      by $\set{W_{n_k}}_{k=1}^{\infty}$ such that $W_{n_{k}} \geq k$
      for all $k \in \N$.
    \item[(2)] By the termwise comparison test we have that
      \[
        \sum_{n=1}^{\infty} \abs{a_n} \le
        \sum_{n=1}^{\infty} \frac{1}{n^2} = \frac{\pi^2}{6},
      \]
      so $\sum_{n=1}^{\infty} a_n$ converges absolutely and in particular
      it converges.
      However we notice that
      \[
        a_n W_n =
        \begin{cases}
          \frac{W_n}{k^2}, &\text{if } n = n_k \text{ for } k=1,2,\dots; \\
          0, &\text{otherwise}.
        \end{cases}
      \]
      But since for $n = n_k$ for $k=1,2,\dots$ we have
      \[
        a_{n_k} W_{n_k} =
        \frac{W_{n_k}}{k^2} \geq \frac{k}{k^2} = \frac{1}{k},
      \]
      we get that
      \[
        \sum_{n=1}^{\infty} a_n W_n =
        \sum_{k=1}^{\infty} a_{n_k} W_{n_k} =
        \sum_{k=1}^{\infty} \frac{1}{k}.
      \]
      Since the harmonic series does not converge $\sum_{n=1}^{\infty} a_n W_n$
      does not converge which completes the proof.
  \end{enumerate}
\end{solution}

\begin{exercise}
  Let $\set{w_n}_{n=1}^{\infty}$ be a sequence of positive numbers satifying
  $w_n \taking{n \to \infty} 0$.
  Define the bounded linear operator $A$ on $\ell^2$ by
  \[
    Av = (w_1 v_1, w_2 v_2, \dots),
  \]
  for $v = (v_1,v_2,\cdots) \in \ell^2$.
  \begin{enumerate}
    \item[(1)] Show that $\norm{A} = \max\limits_{n \geq 1} \abs{w_n}$.
    \item[(2)] Show that $w_n$ is an eigenvalue of $A$ for all $n \in \N$.
    \item[(3)] Show that the point spectrum $\sigma_p(A)$ of $A$ is equal
      to $\set{w_n}_{n=1}^{\infty}$.
    \item[(4)] Show that the spectrum $\sigma(A)$ of $A$ includes $0$.
    \item[(5)] Show that $A$ is injective.
    \item[(6)] In light of $(4)$ and $(5)$, we get $\im(A) \neq \ell^2$.
      Find a vector $v \in \ell^2$ satisfying $v \notin \im(A)$.
  \end{enumerate}
\end{exercise}
\begin{solution} \phantom{}
  \begin{enumerate}
    \item[(1)] Let $v \in \ell^2$ be with norm $1$.
      Then we have
      \[
        \norm{Av} =
        \sqrt{\sum_{n=1}^{\infty} \abs{w_n v_n}^2} \le
      \sqrt{\del{\max_{n \geq 1} w_n}^2 \cdot \sum_{n=1}^{\infty} \abs{v_n}^2} =
        \max_{n \geq 1} w_n.
      \]
      Therefore $\norm{A} \le \max\limits_{n \geq 1} \abs{w_n}$.
      Next we will show that there exists $v$ with norm $1$
      so that $\norm{Av} = \max\limits_{n \geq 1} \abs{w_n}$
      which completes the proof.
      Let $m$ be the index for which $w_m = \max\limits_{n \geq 1} \abs{w_n}$.
      Define
      \[
        v_i :=
        \begin{cases}
          1, &i = m \\
          0, &\text{otherwise}
        \end{cases}
      \]
      and set $v := (v_i)_{n=1}^{\infty} \in \ell^2$.
      We clearly have that $\norm{v} = 1$ and also
      \[
        \norm{Av} =
        \sqrt{\sum_{n=1}^{\infty} \abs{w_n v_n}^2} =
        \sqrt{\max_{n \geq 1} \abs{w_n}^2} =
        \max_{n \geq 1} w_n.
      \]
      This shows that $\norm{A} = \max\limits_{n \geq 1} \abs{w_n}$ and
      completes the proof.
    \item[(2)] We need to show that $w_n$ is an eigenvalue of $A$
      for every $n \in \N$.
      Let $n \in \N$, define
      \[
        v_i :=
        \begin{cases}
          1, &i = n \\
          0, &\text{otherwise}
        \end{cases}
      \]
      and set $v := (v_i)_{n=1}^{\infty} \in \ell^2$.
      We get that
      \[
        Av =
        (w_i v_i)_{n=1}^{\infty} =
        \begin{cases}
          w_n, &i = n \\
          0, &\text{otherwise}
        \end{cases}.
      \]
      Therefore $Av = w_n v$.
      Since the choice of $n$ we arbitrary we get that
      $w_n$ is an eigenvalue of $A$ for all $n \in \N$ which completes the
      proof.
    \item[(3)] In $(2)$
      we saw that $\set{w_n}_{n=1}^{\infty} \subseteq \sigma_p(A)$.
      We will now show the opposite side of the inclusion.
      Let $\lambda \in \sigma_p(A)$.
      Then by definition there exists $0 \neq v \in \ell^2$ so that
      \[
        Av =
        (w_1 v_1, w_2 v_2, \dots) =
        \lambda (v_1, v_2, \dots) =
        (\lambda v_1, \lambda v_2, \dots).
      \]
      Since $v \neq 0$ there exists $n \in \N$ so that $v_n \neq 0$.
      From the previous equation we get that $w_n v_n = \lambda v_n$,
      and since $v_n \neq 0$ we can divide by $v_n$ and get $w_n = \lambda$.
      Therefore $\lambda \in \set{w_n}_{n=1}^{\infty}$.
      This implies that $\sigma_p(A) \subseteq \set{w_n}_{n=1}^{\infty}$.
      We have shown double sided inclusion so
      $\set{w_n}_{n=1}^{\infty} = \sigma_p(A)$ which completes the proof.
    \item[(4)] Showing that $0 \in \sigma(A)$ is equivalent to showing
      that $A$ is not invertible.
      Assume, by way of contradiction, that $A$ has an inverse.
      Then $A^{-1} \in B(\ell^2)$.
      Define the operator $B \colon \ell^2 \to \C^{\infty}$ by
      \[
        Bv :=
        (w_1^{-1} v_1, w_2^{-1} v_2, \dots)
      \]
      for $v = (v_1,v_2,\dots) \in \ell^2$.
      It is well defined since $\set{w_n}_{n=1}^{\infty}$ are positive
      numbers.
      We clearly get that $A^{-1} = B$ (although formally they are different
      since they have different range).
      Therefore we must have $\im(B) \subseteq \ell^2$.
      However, from \Cref{ex:lemma-of-w} there exists
      a sequence $\set{a_n}_{n=1}^{\infty}$ of nonnegative
      numbers such that
      \[
        \sum_{n=1}^{\infty} a_n < \infty \tand
        \sum_{n=1}^{\infty} a_n W_n = \infty.
      \]
      This implies that $v := (a_n)_{n=1}^{\infty} \in \ell^2$
      but $Av \notin \ell^2$.
      Therefore $\im(B) \nsubseteq \ell^2$.
      This contradiction completes the proof.
    \item[(5)] Suppose that $v = (v)_{i=1}^{\infty} \in \ell^2$
      is a vector such that $Av = 0 \in \ell^2$.
      Then
      \[
        Av =
        (w_1 v_1, w_2 v_2, \dots) =
        (0, 0, \dots).
      \]
      This gives $w_n v_n = 0$ for all $n \in \N$.
      Since $\set{w_n}_{n=1}^{\infty}$ are positive we
      can divide by $w_n$ and get $v_n = 0$ for all $n \in \N$
      so $v = 0$ and thus $\ker(A) = 0$
      so $A$ must be injective.
    \item[(6)] Define the sequence $W_n := w_n^{-2}$.
      Let $\set{a_n}_{n=1}^{\infty}$ be the corresponding sequence from
      \Cref{ex:lemma-of-w}.
      Now define $v := (\sqrt{a_n})_{n=1}^{\infty}$.
      Since
      \[
        \norm{v}^2 =
        \sum_{n=1}^{\infty} \sqrt{a_n}^2 =
        \sum_{n=1}^{\infty} a_n < \infty
      \]
      we get that $v \in \ell^2$.
      We claim that $v \notin \im(A)$.
      Assume, by way of contradiction that there exists $u \in \ell^2$
      so that $Au = v$.
      Then
      \[
        Au =
        (w_n u_n)_{n=1}^{\infty} =
        (\sqrt{a_n})_{n=1}^{\infty}.
      \]
      Therefore we must have $u_n = \frac{\sqrt{a_n}}{w_n}$.
      This implies that
      \[
        \norm{u}^2 =
        \sum_{n=1}^{\infty} \del{\frac{\sqrt{a_n}}{w_n}}^2 =
        \sum_{n=1}^{\infty} a_n W_n = \infty.
      \]
      Therefore $u \notin \ell^2$.
      This contradiction completes the proof.
  \end{enumerate}
\end{solution}

\newpage

\section{Compact operators and the Fourier transform}

\begin{exercise}
  In this exercise we will solve the equation
  \[
    \int_{-\infty}^{\infty} f(t) f(x - t) \dt =
    \frac{1}{x^2 + 1}.
  \]
\end{exercise}
\begin{solution}
  We notice that the LHS of the equation is simply $f * f$.
  From the convolution theorem we know
  that $\widehat{f * f} = \hat{f} \cdot \hat{f}$.
  We have also previously seen that if $g(x) = \frac{1}{x^2 + 1}$
  its Fourier transform is
  \[
    \hat{g}(\omega) = \pi e^{-\abs{\omega}}.
  \]
  Therefore
  \[
    \hat{f} \cdot \hat{f} = \pi e^{-\abs{\omega}}
    \implies
    \hat{f}(\omega) = \pm \sqrt{\pi} e^{-\frac{\abs{\omega}}{2}}.
  \]
  We can now notice that $\hat{f}$ is simply $\pm \pi^{-1/2} \hat{g}(\omega/2)$.
  Since the Fourier transformation $\mathcal F_1$ is linear
  and from the
  identity $\mathcal F_1[g(ax)](\omega) = a^{-1} \mathcal F_1[g(x)](\omega/a)$
  we get that
  \[
    \mathcal F_1 \left[\frac{1}{\sqrt{\pi}} \cdot \frac{2}{(2x)^2 + 1}\right](\omega) =
    \frac{2}{\sqrt{\pi}} \mathcal F_1 \left[\frac{1}{(2x)^2 + 1}\right](\omega) =
    \frac{1}{\sqrt{\pi}} \mathcal F_1 \left[\frac{1}{x^2 + 1}\right](\frac{\omega}{2}) =
    \sqrt{\pi} e^{-\frac{\abs{\omega}}{2}}.
  \]
  And similarly we get
  \[
    \mathcal F_1 \left[-\frac{1}{\sqrt{\pi}} \cdot \frac{2}{(2x)^2 + 1}\right](\omega) =
    -\sqrt{\pi} e^{-\frac{\abs{\omega}}{2}}.
  \]
  Therefore the functions
  \begin{gather*}
    f_1 = \frac{1}{\sqrt{\pi}} \cdot \frac{2}{4x^2 + 1} \\
    f_2 = -\frac{1}{\sqrt{\pi}} \cdot \frac{2}{4x^2 + 1}
  \end{gather*}
  are both solutions to the equation.
  Since the Fourier transform $\mathcal F_1$ is injective,
  and we must
  have $\mathcal F_1[f](\omega) = \sqrt{\pi} e^{-\frac{\abs{\omega}}{2}}$,
  these are all the solutions of the equation.
\end{solution}

\begin{exercise}
  Let $a,b > 0$.
  Calculate $\int_{0}^{\infty} \frac{1}{(x^2 + a^2) (x^2 + b^2)} \dx$.
\end{exercise}
\begin{solution}
  The integrand is an even function so we have
  \[
    \int_{0}^{\infty} \frac{1}{(x^2 + a^2) (x^2 + b^2)} \dx =
    \half \int_{-\infty}^{\infty} \frac{1}{(x^2 + a^2) (x^2 + b^2)} \dx.
  \]
  Recall that we saw in the recitations that
  \[
    \mathcal F_1\left[e^{-a |x|}\right](\omega) =
    \frac{2a}{a^2 + \omega^2}.
  \]
  We can now define $f_{y}(x) := e^{-y |x|}$ and see that
  from the convolution theorem
  \[
    \mathcal F_1[f_a * f_b](\omega) =
    \mathcal F_1[f_a](\omega) \cdot \mathcal F_1[f_b](\omega) =
    \frac{4ab}{(a^2 + \omega^2) (b^2 + \omega^2)}.
  \]
  A tiresome computation shows that $f_a * f_b \in C^1$
  and since $\widehat{f_a * f_b} \in \mathcal L^1(\R)$
  we can use the inverse Fourier transform at $x = 0$ to get
  \[
    (f_a * f_b)(0) =
    \frac{1}{2\pi}
    \int_{-\infty}^{\infty} \frac{4ab}{(a^2 + \omega^2) (b^2 + \omega^2)} \cdot
    e^{i \omega \cdot 0} \domega =
    \frac{4ab}{\pi}
    \int_{0}^{\infty} \frac{1}{(a^2 + \omega^2) (b^2 + \omega^2)} \domega.
  \]
  We can compute directly and get,
  \[
    (f_a * f_b)(0) =
    \int_{-\infty}^{\infty} e^{-(a + b)|x|} \dx =
    2 \int_{0}^{\infty} e^{-(a + b) x} \dx =
    \frac{2}{a + b}.
  \]
  By plugging this computation in our previous equality we can easily conclude
  that
  \[
    \int_{0}^{\infty} \frac{1}{(x^2 + a^2) (x^2 + b^2)} \dx =
    \frac{\pi}{2ab (a + b)}.
  \]
\end{solution}

\begin{exercise}
  Let $f(x) = \ind_{[-b,b]}$.
  We have previously seen
  that $\hat{f}(\omega) = \frac{2 \sin(b \omega)}{\omega}$.
  \begin{enumerate}
    \item[(1)] Show that
      \[
        (f * f)(x) =
        \begin{cases}
          0, &|x| > 2b; \\
          2b - x, &0 \le x \le 2b; \\
          2b + x, &-2b \le x \le 0.
        \end{cases}
      \]
    \item[(2)] Find the Fourier transform of $(f * f)$.
    \item[(3)] Use parts $(1)$ and $(2)$
      to calculate the following integral:
      \[
        \int_{0}^{\infty} \frac{\sin^4(b \omega)}{\omega^4} \domega.
      \]
  \end{enumerate}
\end{exercise}
\begin{solution} \phantom{}
\begin{enumerate}
  \item[(1)] By definition
    \[
      (f * f)(x) =
      \int_{-\infty}^{\infty} \ind_{[-b,b]} (t) \cdot \ind_{[-b,b]} (x - t) \dt.
    \]
    We know that $\ind_{[-b,b]} (x - t) = \ind_{[x-b,x+b](t)}$,
    thus we can write
    \[
      (f * f)(x) =
      \int_{-\infty}^{\infty} \ind_{[-b,b] \cap [x-b,x+b]} (t) \dt =
      \int_{[-b,b] \cap [x-b,x+b]} \dt.
    \]
    Whenever $|x| > 2b$ we get $[-b,b] \cap [x-b,x+b] = \emptyset$
    so $(f * f)(x) = 0$.
    When $0 \le x \le 2b$ the intersection between the intervals
    is $[x - b, b]$, and when $-2b \le x \le 0$ the intersection is
    $[-b, x + b]$.
    In both of these cases the integral is immediate and equals
    $2b - x$ and $2b + x$ respectively.
    Therefore
    \[
        (f * f)(x) =
        \begin{cases}
          0, &|x| > 2b; \\
          2b - x, &0 \le x \le 2b; \\
          2b + x, &-2b \le x \le 0.
        \end{cases}
      \]
    which completes the proof.
  \item[(2)] We need to find $\mathcal F_1[f * f](\omega)$.
    Since $f \in \mathcal L^1(\R)$ we are allowed to use the convlusion
    theorem which, coupled with the identity
    from the recitations $\hat{f}(\omega) = \frac{2 \sin(b \omega)}(\omega)$,
    gives the result immediately:
    \[
      \mathcal F_1[f * f](\omega) =
      \hat{f}(\omega) \cdot \hat{f}(\omega) =
      \frac{4 \sin^2(b \omega)}{\omega^2}.
    \]
  \item[(3)] Since the integrand is an even function we get that
    \[
      \int_{0}^{\infty} \frac{\sin^4(b \omega)}{\omega^4} \domega =
      \half \int_{-\infty}^{\infty} \frac{\sin^4(b \omega)}{\omega^4} \domega.
    \]
    We can now use the result from $(2)$ to get that
    \[
      \int_{-\infty}^{\infty} \frac{\sin^4(b \omega)}{\omega^4} \domega =
      \int_{-\infty}^{\infty} \abs{\frac{\sin^2(b \omega)}{\omega^2}}^2 \domega =
      \frac{1}{16} \norm{\mathcal F_1[f * f]}_2^2.
    \]
    Since $f \in \mathcal L^2(\R)$ we can use Plancherel's theorem
    to obtain
    that
    $\norm{\mathcal F_1[f * f]}_2 = \sqrt{2 \pi} \cdot \norm{f * f}_2$.
    Therefore
    \[
      \int_{-\infty}^{\infty} \frac{\sin^4(b \omega)}{\omega^4} \domega =
      \frac{\pi}{8} \cdot \norm{f * f}_2^2.
    \]
    We can now use the calculation from $(1)$ to compute the norm:
    \[
      \norm{f * f}_2^2 =
      \int_{-\infty}^{\infty} \abs{(f * f)(x)}^2 \dx = 
      \int_{-2b}^{0} (2b + x)^2 \dx +
      \int_{0}^{2b} (2b - x)^2 \dx.
    \]
    We can find these integrals using a simple substitution and get that
    \[
      \norm{f * f}_2^2 =
      \frac{8b^3}{3} + \frac{8b^3}{3} =
      \frac{16b^3}{3}.
    \]
    Therefore we have that
    \[
      \int_{0}^{\infty} \frac{\sin^4(b \omega)}{\omega^4} \domega =
      \half \int_{-\infty}^{\infty} \frac{\sin^4(b \omega)}{\omega^4} \domega =
      \half \del{\frac{\pi}{8} \cdot \frac{16b^3}{3}} =
      \frac{\pi b^3}{3}.
    \]
  \end{enumerate}
\end{solution}



\end{document}
